
%\needspace{4\baselineskip}
\color{uniblue}{\subsection{Управляющие воздействия на СКРМ (УВ СКРМ)}}
\color{black}

\begin{enumerate}

\item
Для реализации управляющих воздействий от противоаварийной автоматики на СКРМ предусматриваются функции УВ СКРМ. В устройстве предусмотрено две функции УВ СКРМ («УВ СКРМ(n)»), здесь и далее в этом подразделе: (n) -- номер функции, $n\in\{1,2\}$.

\item При срабатывании защит СКРМ прохождение команды на отключение СКРМ от противоаварийной автоматики блокируется, блокировка УВ СКРМ остается в памяти триггера до прихода команды сброса. Сброс триггера осуществляется при помощи ФК «Сброс блок УВ СКРМ (n)».

\item
Ввод УВ СКРМ в работу на этапе параметрирования устройства осуществляется своим программным переключателем SGF1 «Ввод\_функции». Информация о введенной функции отображается выходным сигналом «УВ СКРМ(n): Ввод».
\item
Функция УВ СКРМ может быть выведена из работы оперативно путем активации сигнала «ОВ УВ СКРМ(n)». Выведенное состояние функции характеризуется активным состоянием сигнала «УВ СКРМ(n): Оперативный вывод».

\item
При активном состоянии входного сигнала «УВ СКРМ(n) на сигн» (ключ «УВ СКРМ(n) на сигн» введен) действие УВ СКРМ(n) переводится на сигнал.

\item
Алгоритм работы функции выполнен в соответствии с рисунком \ref{uvskrm:img1}. В таблице \ref{uvskrm:tbl1} приведены параметры, необходимые для настройки функции.

\vspace{3mm}
\begin{figure}[H]
\centering
\includegraphics[width=0.85\textwidth,height=0.85\textheight,keepaspectratio]{\fbpath/208. УВ СКРМ/_latex/img/uvskrm.pdf}
\captionof{figure}{Функциональная схема логики УВ СКРМ}\label{uvskrm:img1}
\end{figure}

\end{enumerate}
